% Options for packages loaded elsewhere
% Options for packages loaded elsewhere
\PassOptionsToPackage{unicode}{hyperref}
\PassOptionsToPackage{hyphens}{url}
\PassOptionsToPackage{dvipsnames,svgnames,x11names}{xcolor}
%
\documentclass[
  ngerman,
  10pt,
  twocolumn,
  twoside]{article}
\usepackage{xcolor}
\usepackage[top=2.5cm,bottom=2.5cm,left=2cm,right=2cm]{geometry}
\usepackage{amsmath,amssymb}
\setcounter{secnumdepth}{5}
\usepackage{iftex}
\ifPDFTeX
  \usepackage[T1]{fontenc}
  \usepackage[utf8]{inputenc}
  \usepackage{textcomp} % provide euro and other symbols
\else % if luatex or xetex
  \usepackage{unicode-math} % this also loads fontspec
  \defaultfontfeatures{Scale=MatchLowercase}
  \defaultfontfeatures[\rmfamily]{Ligatures=TeX,Scale=1}
\fi
\usepackage{lmodern}
\ifPDFTeX\else
  % xetex/luatex font selection
\fi
% Use upquote if available, for straight quotes in verbatim environments
\IfFileExists{upquote.sty}{\usepackage{upquote}}{}
\IfFileExists{microtype.sty}{% use microtype if available
  \usepackage[]{microtype}
  \UseMicrotypeSet[protrusion]{basicmath} % disable protrusion for tt fonts
}{}
\usepackage{setspace}
\makeatletter
\@ifundefined{KOMAClassName}{% if non-KOMA class
  \IfFileExists{parskip.sty}{%
    \usepackage{parskip}
  }{% else
    \setlength{\parindent}{0pt}
    \setlength{\parskip}{6pt plus 2pt minus 1pt}}
}{% if KOMA class
  \KOMAoptions{parskip=half}}
\makeatother
% Make \paragraph and \subparagraph free-standing
\makeatletter
\ifx\paragraph\undefined\else
  \let\oldparagraph\paragraph
  \renewcommand{\paragraph}{
    \@ifstar
      \xxxParagraphStar
      \xxxParagraphNoStar
  }
  \newcommand{\xxxParagraphStar}[1]{\oldparagraph*{#1}\mbox{}}
  \newcommand{\xxxParagraphNoStar}[1]{\oldparagraph{#1}\mbox{}}
\fi
\ifx\subparagraph\undefined\else
  \let\oldsubparagraph\subparagraph
  \renewcommand{\subparagraph}{
    \@ifstar
      \xxxSubParagraphStar
      \xxxSubParagraphNoStar
  }
  \newcommand{\xxxSubParagraphStar}[1]{\oldsubparagraph*{#1}\mbox{}}
  \newcommand{\xxxSubParagraphNoStar}[1]{\oldsubparagraph{#1}\mbox{}}
\fi
\makeatother


\usepackage{longtable,booktabs,array}
\usepackage{calc} % for calculating minipage widths
% Correct order of tables after \paragraph or \subparagraph
\usepackage{etoolbox}
\makeatletter
\patchcmd\longtable{\par}{\if@noskipsec\mbox{}\fi\par}{}{}
\makeatother
% Allow footnotes in longtable head/foot
\IfFileExists{footnotehyper.sty}{\usepackage{footnotehyper}}{\usepackage{footnote}}
\makesavenoteenv{longtable}
\usepackage{graphicx}
\makeatletter
\newsavebox\pandoc@box
\newcommand*\pandocbounded[1]{% scales image to fit in text height/width
  \sbox\pandoc@box{#1}%
  \Gscale@div\@tempa{\textheight}{\dimexpr\ht\pandoc@box+\dp\pandoc@box\relax}%
  \Gscale@div\@tempb{\linewidth}{\wd\pandoc@box}%
  \ifdim\@tempb\p@<\@tempa\p@\let\@tempa\@tempb\fi% select the smaller of both
  \ifdim\@tempa\p@<\p@\scalebox{\@tempa}{\usebox\pandoc@box}%
  \else\usebox{\pandoc@box}%
  \fi%
}
% Set default figure placement to htbp
\def\fps@figure{htbp}
\makeatother



\ifLuaTeX
\usepackage[bidi=basic]{babel}
\else
\usepackage[bidi=default]{babel}
\fi
% get rid of language-specific shorthands (see #6817):
\let\LanguageShortHands\languageshorthands
\def\languageshorthands#1{}
\ifLuaTeX
  \usepackage[german]{selnolig} % disable illegal ligatures
\fi


\setlength{\emergencystretch}{3em} % prevent overfull lines

\providecommand{\tightlist}{%
  \setlength{\itemsep}{0pt}\setlength{\parskip}{0pt}}



 
\usepackage[style=authoryear,backend=biber,style=authoryear-comp,uniquename=false,uniquelist=false]{biblatex}
\addbibresource{references.bib}


\usepackage{csquotes}
\usepackage{booktabs}
\usepackage{array}
\usepackage{multirow}
\usepackage{wrapfig}
\usepackage{float}
\usepackage{colortbl}
\usepackage{pdflscape}
\usepackage{threeparttable}
\usepackage{threeparttablex}
\usepackage[normalem]{ulem}
\usepackage{makecell}
\usepackage{xcolor}
\usepackage{multicol}
\usepackage{fancyhdr}
\usepackage{titlesec}
\usepackage{etoolbox}
\usepackage{tabularx}
\usepackage{afterpage}
\usepackage{tikz}

% Suppress default maketitle
\renewcommand{\maketitle}{}

% Define WISTA colors (Havelock Blue as main)
\definecolor{wistaBlue}{RGB}{85,145,200}
\definecolor{wistaLightBlue}{RGB}{120,170,215}
\definecolor{wistaDarkGray}{RGB}{64,64,64}
\definecolor{wistaLightGray}{RGB}{128,128,128}

% Custom section formatting - number on top with column-width underline
\titleformat{\section}[display]
  {\normalfont\large\bfseries}
  {\color{wistaBlue}\large\bfseries\thesection\vspace{0.1cm}\\\noindent\rule{\columnwidth}{1pt}}
  {0.3cm}
  {\color{wistaDarkGray}\large\bfseries}
  [\vspace{0.1cm}\noindent{\color{wistaBlue}\rule{\columnwidth}{0.5pt}}\vspace{0.3cm}]
\titlespacing*{\section}{0pt}{2ex plus 1ex minus .2ex}{1ex plus .2ex}

% Subsection formatting - number and title with underline
\titleformat{\subsection}
  {\color{wistaDarkGray}\normalsize\bfseries}
  {\color{wistaBlue}\normalsize\bfseries\thesubsection}
  {0.5em}
  {\color{wistaDarkGray}\normalsize\bfseries}
  [\vspace{0.05cm}\noindent{\color{wistaBlue}\rule{0.5\columnwidth}{0.4pt}}\vspace{0.1cm}]
  
% Subsubsection formatting
\titleformat{\subsubsection}
  {\color{wistaDarkGray}\small\bfseries}
  {\color{wistaBlue}\small\bfseries\thesubsubsection}
  {0.3em}
  {\color{wistaDarkGray}\small\bfseries}

% Paragraph formatting (4th level)
\titleformat{\paragraph}[runin]
  {\color{wistaDarkGray}\normalsize\bfseries}
  {}
  {0em}
  {}
  [.]
\titlespacing*{\paragraph}{0pt}{1ex plus 0.5ex minus 0.2ex}{0.5em}

% Header and footer setup with blue line under header
\pagestyle{fancy}
\fancyhf{}
\fancyhead[L]{\small\textit{Oliver Hauke, Annette Kristiansen}}
\fancyhead[R]{\small\textit{Effiziente Analyse großer Forschungsdaten}}
\fancyfoot[L]{\small\thepage}
\fancyfoot[R]{\small Statistisches Bundesamt | WISTA | {\color{wistaBlue}6} | 2025}
\renewcommand{\headrulewidth}{0pt}
\renewcommand{\footrulewidth}{0pt}
% Add blue line under header
\renewcommand{\headrule}{\color{wistaBlue}\hrule width\headwidth height 1pt \vskip 2pt}
\renewcommand{\headrulewidth}{1pt}

% Make all figures and tables span both columns (full width)
\let\origfigure\figure
\let\endorigfigure\endfigure
\renewenvironment{figure}[1][htbp]{%
  \origfigure*[#1]%
}{%
  \endorigfigure*%
}

\let\origtable\table
\let\endorigtable\endtable
\renewenvironment{table}[1][htbp]{%
  \origtable*[#1]%
}{%
  \endorigtable*%
}

% Footnotes stay within column
\usepackage[stable]{footmisc}

% Custom bullet points with ">"
\usepackage{enumitem}
\setlist[itemize]{label={\color{wistaBlue}\textgreater}, leftmargin=1.5em, itemsep=0pt, parsep=0pt}

% Blue arrow for cross-references (bottom-right arrow)
\usepackage{amssymb}

% Redefine ref to include blue arrow for figure/table/section references
\let\oldref\ref
\renewcommand{\ref}[1]{{\color{wistaBlue}$\searrow$}\oldref{#1}}
\makeatletter
\@ifpackageloaded{caption}{}{\usepackage{caption}}
\AtBeginDocument{%
\ifdefined\contentsname
  \renewcommand*\contentsname{Inhaltsverzeichnis}
\else
  \newcommand\contentsname{Inhaltsverzeichnis}
\fi
\ifdefined\listfigurename
  \renewcommand*\listfigurename{Abbildungsverzeichnis}
\else
  \newcommand\listfigurename{Abbildungsverzeichnis}
\fi
\ifdefined\listtablename
  \renewcommand*\listtablename{Tabellenverzeichnis}
\else
  \newcommand\listtablename{Tabellenverzeichnis}
\fi
\ifdefined\figurename
  \renewcommand*\figurename{Abbildung}
\else
  \newcommand\figurename{Abbildung}
\fi
\ifdefined\tablename
  \renewcommand*\tablename{Tabelle}
\else
  \newcommand\tablename{Tabelle}
\fi
}
\@ifpackageloaded{float}{}{\usepackage{float}}
\floatstyle{ruled}
\@ifundefined{c@chapter}{\newfloat{codelisting}{h}{lop}}{\newfloat{codelisting}{h}{lop}[chapter]}
\floatname{codelisting}{Listing}
\newcommand*\listoflistings{\listof{codelisting}{Listingverzeichnis}}
\makeatother
\makeatletter
\makeatother
\makeatletter
\@ifpackageloaded{caption}{}{\usepackage{caption}}
\@ifpackageloaded{subcaption}{}{\usepackage{subcaption}}
\makeatother
\usepackage{bookmark}
\IfFileExists{xurl.sty}{\usepackage{xurl}}{} % add URL line breaks if available
\urlstyle{same}
\hypersetup{
  pdftitle={Effiziente Analyse großer Forschungsdaten mithilfe von R},
  pdfauthor={Oliver Hauke; Annette Kristiansen},
  pdflang={de},
  colorlinks=true,
  linkcolor={blue},
  filecolor={Maroon},
  citecolor={Blue},
  urlcolor={Blue},
  pdfcreator={LaTeX via pandoc}}


\title{Effiziente Analyse großer Forschungsdaten mithilfe von R}
\author{Oliver Hauke \and Annette Kristiansen}
\date{2025-11-23}
\begin{document}
\maketitle

% Custom WISTA-style title page
\begin{titlepage}
\thispagestyle{empty}
\pagestyle{empty}
\onecolumn

% Horizontal blue line at top (header line)
\noindent{\color{wistaBlue}\rule{\textwidth}{1pt}}

\vspace{-0.5cm}

% Vertical blue line on the left (6.3cm from left edge, starts below header line, stops above footer)
\begin{tikzpicture}[remember picture, overlay]
  \draw[wistaBlue, line width=2pt] 
    ([xshift=6.3cm, yshift=-4cm]current page.north west) -- 
    ([xshift=6.3cm, yshift=3.5cm]current page.south west);
\end{tikzpicture}

\vspace*{0.5cm}

% Left side - short author bios (right-aligned to the vertical line)
\begin{minipage}[t]{3.9cm}
\selectlanguage{ngerman}%
\raggedleft
\hyphenpenalty=0
\tolerance=1000
\fontsize{7}{8.5}\selectfont
\vspace{0pt}%
{\color{wistaBlue}\textbf{Oliver Hauke}}\\[0.15cm]
ist Mathematiker und IT-Referent im Kompetenzzentrum „Analyse und Auswertung" des Statistischen Bundesamtes. Er betreut die technische Infrastruktur des Forschungsdatenzentrums im Statistischen Bundesamt und Arbeiten für das Netzwerk empirische Steuerforschung.\\[0.5cm]

{\color{wistaBlue}\textbf{Annette Kristiansen}}\\[0.15cm]
ist Ökonomin und Referentin im Referat „Unternehmenssteuern" des Statistischen Bundesamtes. Sie beschäftigt sich mit der Zusammenführung der Unternehmenssteuerstatistiken und betreut das Netzwerk empirische Steuerforschung. Für ihre externe Promotion an der Universität Bremen untersucht sie die technologische Vielfalt multinationaler Unternehmen.

\end{minipage}%
\hspace{0.4cm}% Gap to vertical line (left side)
\hspace{0.4cm}% Gap from vertical line (right side)
% Right side - title, authors, keywords, abstracts
\begin{minipage}[t]{11.8cm}
\raggedright
\vspace{1.2cm}%

% Title (uppercase via command, not bold)
{\LARGE\color{wistaDarkGray}\begin{spacing}{1.3}\MakeUppercase{Effiziente Analyse großer Forschungsdaten mithilfe von R}\end{spacing}}

\vspace{0.15cm}

% Blue horizontal line (underline for title)
{\color{wistaBlue}\rule{\linewidth}{0.5pt}}

\vspace{0.5cm}

% Authors inline, dark gray
{\large\color{wistaDarkGray}\bfseries Oliver Hauke, Annette Kristiansen}

%\vspace{0.15cm}

% Blue line (underline for authors)
{\color{wistaBlue}\rule{\linewidth}{0.5pt}}

\vspace{0.5cm}

% Keywords (blue arrow then bold label and blue text)
{\color{wistaBlue}\footnotesize$\searrow$\normalsize\ {\bfseries Schlüsselwörter:} Effiziente Analysen, R-Programmierung, Forschungsdaten, Unternehmenssteuerstatistik, Business-Tax-Panel, Netzwerk empirische Steuerforschung}

\vspace{0.3cm}

% Zusammenfassung (uppercase via command)
{\color{wistaDarkGray}\bfseries\MakeUppercase{Zusammenfassung}}\\[0.15cm]
\small
Das Forschungsdatenzentrum des Statistischen Bundesamtes baut seine technische Infrastruktur gezielt aus, um der Wissenschaft erweiterte Analysemöglichkeiten amtlicher Daten bereitzustellen. Ab November 2025 steht Forschenden exklusiver Zugang zu diesen erweiterten Systemressourcen in R und Stata zur Verfügung. R-basierte Tools -- insbesondere Apache Arrow, Parquet und DuckDB -- ermöglichen es, gängige Analysen großer Forschungsdatensätze in Echtzeit auszuführen. Am Beispiel des Business-Tax-Panel (2013 bis 2019) werden die erzielbaren Effizienzgewinne bei der Auswertung umfangreicher steuerstatistischer Daten im Rahmen des Netzwerks empirische Steuerforschung demonstriert.

\vspace{0.5cm}

% English keywords and abstract (italicized)
{\itshape\color{wistaBlue}\footnotesize$\searrow$\normalsize\ {\bfseries Keywords:} Efficient analysis -- R programming -- research data -- corporate tax statistics -- Business Tax Panel -- empirical tax research network}

\vspace{0.3cm}

{\itshape\color{wistaDarkGray}\bfseries\MakeUppercase{Abstract}}\\[0.15cm]
\small\itshape
The Research Data Centre at the Federal Statistical Office is expanding its technical infrastructure to provide researchers with enhanced analytical capabilities for official data. From November 2025, researchers will have exclusive access to these extended system resources in R and Stata. R-based tools -- particularly Apache Arrow, Parquet, and DuckDB -- enable real-time analysis of large research datasets. Using the Business Tax Panel (2013--2019), the achievable efficiency gains in analysing large-scale tax statistics within the empirical tax research network are demonstrated.

\end{minipage}

\vfill

\vspace{0.5cm}
% Footer matching document style
\noindent\makebox[\textwidth]{%
  \small 1%
  \hfill%
  \small Statistisches Bundesamt | WISTA | {\color{wistaBlue}6} | 2025%
}

\end{titlepage}

% Switch back to two-column mode and restore fancy page style
\twocolumn
\pagestyle{fancy}


\setstretch{1}



\printbibliography



\end{document}
